\chapter{Evaluation}\label{s:evaluation}

This chapter presents the experimental setup, evaluation runs, and results
for all four research questions. Interpretation of the results and their
implications are deferred to Chapter~\ref{s:discussion}.

% ──────────────────────────────────────────────────────────────────────────────
\section{Experimental Setup}\label{s:eval-setup}

\textbf{Dataset.} The evaluation dataset consists of 584 unique
intent--BraneScript pairs, authored manually across nine source files. A
deduplication pass reduced an initial corpus of 658 examples to 584 by
removing redundant entries, which were then split into a training set of 497
examples and a held-out test set of 87 examples. The test set was not used
during any stage of training or system development. The dataset covers all
seven Brane packages used in this project (\texttt{healthcare},
\texttt{genomics}, \texttt{epidemics}, \texttt{statistics},
\texttt{text\_analysis}, \texttt{data\_masking}, \texttt{datetime}), with
intent complexity ranging from single-package, single-function invocations to
multi-package workflows with intermediate results passed across five or more
steps. An additional 1974 paraphrases (three per intent, generated by
Qwen3.5-4B) were produced for the semantic cache evaluation
(Section~\ref{s:eval-cache}).

\textbf{Models.} Three base models are evaluated as zero-shot baselines:
Qwen3.5-4B, Qwen3.5-9B, and Qwen3.6-27B, all from the Qwen3 family to enable
clean attribution of performance differences to model scale. Two supervised
fine-tuning runs are also evaluated: a 9B model fine-tuned for three epochs on
the training set with the original system prompt (\texttt{qwen3.5-9b-3ep}),
and a 9B model fine-tuned for five epochs on the training set with the
corrected system prompt (\texttt{qwen3.5-9b-ep5}).

\textbf{Pipeline.} Unless stated otherwise, every evaluation run uses the
complete three-stage pipeline: (1) the \emph{intent decomposer} reformulates
the user's natural-language intent into a structured set of sub-tasks with
explicit package and function candidates; (2) the \emph{RAG retriever} fetches
the most relevant package documentation and example BraneScript fragments per
sub-task from the retrieval index; and (3) the \emph{generator} produces
a BraneScript workflow conditioned on the decomposed intent and retrieved
context. The decomposer ablation in Section~\ref{s:eval-decomp} is the only
run that deviates from this configuration.

\textbf{Metrics.} The three primary metrics used throughout this chapter ---
compile rate, execution rate, and output match rate --- are defined in
Section~\ref{s:eval-metrics}. The cache-specific metrics (hit rate, correct
rate, and false-positive rate) are defined in
Section~\ref{s:cache-metrics}.

\textbf{Hardware.} All generation and fine-tuning runs execute on Snellius,
the Dutch national HPC facility, on H100 GPU nodes allocated via SLURM.
BraneScript execution and metric collection run on a local workstation with
Docker-deployed Brane services.

% ──────────────────────────────────────────────────────────────────────────────
\section{RQ1: Translation Accuracy Across Model Scales}\label{s:eval-rq1}

\subsection{Main results}

Figure~\ref{fig:base-model-comparison} shows compile rate, execution rate, and
output match rate for all three base models evaluated against the 497 test
intents.

\begin{figure}[htbp]
    \centering
    \includegraphics[width=\textwidth]{base_model_comparison.png}
    \caption{Compile rate, execution rate, and output match rate for the three
    base models. Each metric group is shown as a cluster of three bars, one
    per model ($n = 497$).}
    \label{fig:base-model-comparison}
\end{figure}

All three metrics improve monotonically with model scale. The 4B model achieves
a compile rate of 56.4\%, an execution rate of 54.6\%, and an output match
rate of 24.1\%. At 9B parameters, compile rate rises to 66.1\% and output
match to 36.9\%. The largest jump occurs at the 9B-to-27B transition: compile
rate reaches 95.6\%---a gain of 29.5 percentage points---while execution rate
follows closely at 93.2\%. Output match at 27B stands at 49.6\%, an
improvement of 12.7 percentage points over the 9B model.

The gap between execution rate and output match rate widens with scale. For the
4B model the gap is 30.5 percentage points; for the 27B model it is 43.6
percentage points. This indicates that as model scale increases, the dominant
remaining failure mode shifts from the script failing to compile or execute to
the script producing a syntactically and semantically valid workflow whose
output does not exactly match the reference.

\subsection{Error type breakdown}\label{s:eval-errors}

To understand \emph{where} each model fails rather than merely \emph{how
often}, every failed evaluation was classified into one of five categories:

\begin{itemize}
    \item \emph{Compile: syntax error} --- a parse or structural error
          detectable by the BraneScript grammar alone.
    \item \emph{Compile: undefined variable} --- a reference to a variable
          that was not previously bound in the script's scope.
    \item \emph{Compile: undefined function} --- an invocation of a function
          not exported by the imported packages.
    \item \emph{Runtime error} --- a script that compiled successfully but
          failed during Brane package execution.
    \item \emph{Output mismatch} --- a script that executed without error but
          produced output that did not match the reference.
\end{itemize}

Figure~\ref{fig:error-breakdown} shows the distribution of these
categories for the three base models.

\begin{figure}[htbp]
    \centering
    \includegraphics[width=\textwidth]{error_breakdown.png}
    \caption{Error type breakdown for the three base models. Each bar shows
    the fraction of all 497 test examples in each outcome category. Shades of
    red distinguish the four compile-error sub-types; amber indicates output
    mismatch; green indicates correct output.}
    \label{fig:error-breakdown}
\end{figure}

The 4B model shows the most compile-time failures overall. Undefined variable
errors account for 24.1\% of all 497 test examples, making it the single
largest failure category. Syntax errors add a further 17.5\%, and output
mismatch accounts for 30.5\%. Only 24.1\% of scripts produce a correct result.
Runtime errors and undefined function errors are minor contributors at 1.8\%
and 1.4\% respectively.

The 9B model's profile shifts noticeably: undefined variable errors persist
at 24.5\% but syntax errors drop sharply to 4.8\%, indicating that the larger
model handles basic BraneScript grammar more reliably. Undefined function
errors rise slightly to 3.8\%. Output mismatch decreases to 28.1\% and the
correct fraction rises to 36.9\%.

The 27B model largely eliminates compile failures. Syntax errors fall to 1.4\%,
undefined variable errors to 0.2\%, and undefined function errors to 2.4\%.
The dominant failure category becomes output mismatch at 43.6\%: the model
consistently produces scripts that compile and execute, but whose output
diverges from the reference. Runtime errors account for 2.4\%.

% ──────────────────────────────────────────────────────────────────────────────
\section{RQ2: Effect of Retrieval, Prompt Engineering, and Decomposition}
\label{s:eval-rq2}

\subsection{Two evaluation rounds}\label{s:eval-rounds}

An initial evaluation run was conducted with an early version of the
system prompt and retrieval configuration. Several engineering issues were
subsequently identified and corrected:

\begin{itemize}
    \item The system prompt permitted empty-string function arguments and
          under-specified \texttt{commit\_result} usage.
    \item RAG package aliases were insufficiently precise, causing the retriever
          to return documentation for the wrong package. Package keywords were
          also added to the retrieval index to improve matching precision.
    \item The post-processing step that strips stray markdown tokens from
          generation output was incomplete, causing otherwise valid scripts to
          fail at the BraneScript lexer. Additionally, a repetition penalty
          was added to generation to suppress degenerate repetition loops in
          which the model emitted sequences of backtick characters after
          otherwise valid BraneScript.
    \item The evaluation dataset was audited and intent--BraneScript
          mismatches were corrected across all examples.
\end{itemize}

The same model weights were re-evaluated after all three fixes.
Figure~\ref{fig:baseline-first-run} shows a grouped bar chart with one pair
of bars per model: the left bar (slate blue) represents the pre-fix output
match rate, and the right bar (teal) represents the post-fix rate. For every
model, the teal bar is substantially taller than the slate blue bar, indicating
a consistent improvement across the board. The gap is largest for the 9B model
and smallest for the 4B model in absolute terms, while the 27B model achieves
the highest post-fix rate of the three. Across all three models the output
match rate increases substantially: from 3.6\% to 24.1\% for the 4B model,
from 6.0\% to 36.9\% for the 9B model, and from 13.5\% to 49.6\% for the 27B
model. The relative ordering of the three models is preserved between runs,
with the 27B model consistently outperforming the smaller models on this
metric.

\begin{figure}[H]
    \centering
    \includegraphics[width=0.82\textwidth]{baseline_first_run.png}
    \caption{Output match rate before (Run~1, pre-fix) and after (Run~2,
    post-fix) the three engineering fixes ($n=497$ in both runs). Identical
    model weights were used in both runs.}
    \label{fig:baseline-first-run}
\end{figure}

The 27B model's execution rate rises from 37.9\% to 93.2\% between the two
runs using identical model weights. The 9B model's execution rate rises from
19.5\% to 64.9\%, and the 4B model's from 6.2\% to 54.6\%.

\subsection{Intent decomposition ablation}\label{s:eval-decomp}

The intent decomposer is the first stage of the pipeline. Given a raw
natural-language intent, it reformulates the request into a structured set of
sub-tasks, each annotated with explicit package and function candidates. This
structured representation is then passed to the RAG retriever and the
generator. To isolate the contribution of this stage, an ablation run was
conducted in which the decomposer was entirely disabled: the raw intent string
was passed directly to the retriever and generator without any type of reformulation while keeping the other pipeline components (the retrieval indices, the generator model,
the system prompt, and the post-processing step) identical to the full pipeline configuration. The experiment used Qwen3.5-9B as the generator
and the full 497-example test set.

Figure~\ref{fig:decomp-ablation} presents the results as a grouped bar chart
with two groups on the x-axis: full pipeline as the first one, and pipeline with no decomposition step as the second.
For each run, the compile rate (slate blue), execution
rate (teal), and output match rate (purple) is evaluated. Comparing the two groups, all three bars are
shorter in the no-decomposition group, indicating a consistent drop across every metric when the decomposer is removed.

\begin{figure}[htbp]
    \centering
    \includegraphics[width=0.75\textwidth]{decomp_ablation.png}
    \caption{Effect of disabling the intent decomposition stage on
    Qwen3.5-9B ($n = 497$). Slate blue = compile rate; teal = execution rate;
    purple = output match rate.}
    \label{fig:decomp-ablation}
\end{figure}

Compile rate falls from 66.1\% to 52.9\%, a drop of 13.2 percentage points. Execution rate falls from
64.9\% to 48.1\%, a drop of 16.8 percentage points. Output match rate falls
from 36.9\% to 26.0\%, a drop of 10.9 percentage points. The degradation is
consistent across all three metrics, with execution rate showing the largest
absolute drop and output match the smallest.

% ──────────────────────────────────────────────────────────────────────────────
\section{RQ3: Effect of Supervised Fine-Tuning}\label{s:eval-sft}

\subsection{Fine-tuning configurations and results}

Figure~\ref{fig:sft-results} reports results for both SFT models alongside the
9B base model. \texttt{qwen3.5-9b-3ep} was fine-tuned for three epochs on the
pre-deduplication version of the training set (560 examples) using the
original system prompt. \texttt{qwen3.5-9b-ep5} was fine-tuned for five
epochs on the final deduplicated training set (497 examples) using the
corrected system prompt. Both runs share the following hyperparameters:

\begin{itemize}
    \item \textbf{Method:} QLoRA (quantised LoRA, 4-bit base model)
    \item \textbf{LoRA rank:} 16; \textbf{LoRA $\alpha$:} 32
    \item \textbf{Learning rate:} $2 \times 10^{-4}$ with cosine scheduler
    \item \textbf{Batch size:} 2 per device, gradient accumulation 8 steps (effective batch 16)
    \item \textbf{Max sequence length:} 2048 tokens
    \item \textbf{Optimiser:} AdamW with default betas
\end{itemize}

\begin{figure}[htbp]
    \centering
    \includegraphics[width=0.85\textwidth]{sft_results.png}
    \caption{Compile, execution, and output match rates for the Qwen3.5-9B
    base model and both SFT variants ($n = 497$ for all runs).}
    \label{fig:sft-results}
\end{figure}

Compared to the 9B base, the 3ep SFT model shows a decline across all three
metrics: compile rate drops from 66.1\% to 56.0\%, execution rate from 64.9\%
to 49.4\%, and output match from 36.9\% to 15.9\%. Within the compile failure
categories, syntax errors rise from 4.8\% to 28.7\% and undefined function
errors from 3.8\% to 13.1\%, while undefined variable errors fall sharply from
24.5\% to 0.4\%.

The 5ep SFT model shows a near-total collapse in performance: compile rate
falls to 2.4\%, execution rate to 2.2\%, and output match to 0.8\%. Inspection
of the generated outputs reveals that the majority of scripts are truncated
mid-generation, a consequence of the 512-token output limit that was active
during this evaluation run. These results are therefore not representative of
the model's learned behaviour and are included for completeness only.

% ──────────────────────────────────────────────────────────────────────────────
\section{RQ4: Semantic Cache as a Reproducibility Layer}\label{s:eval-cache}

\textbf{Experiment setup.}
The semantic cache stores previously generated BraneScript workflows indexed
by the embedding of the intent that produced them. When a new intent arrives,
its embedding is compared against all cached entries using cosine similarity.
If the similarity to any cached entry exceeds a configurable threshold, the
stored BraneScript is returned directly without invoking the LLM pipeline.
This allows identical or near-identical intents to bypass generation entirely,
guaranteeing reproducibility for repeated requests.

To evaluate the cache, it was seeded with all 584 training intents paired with
their reference BraneScript. The 1974 paraphrase intents described in
Section~\ref{s:eval-setup} were then used as queries. A paraphrase query is a
semantically equivalent restatement of an original intent using different
surface wording. Since each paraphrase corresponds to a known cached entry,
the expected behaviour is a cache hit returning the correct BraneScript.

Three metrics are reported across eight similarity thresholds from 0.80 to
0.98, as defined in Section~\ref{s:cache-metrics}: hit rate, correct rate,
and false-positive rate.

Table~\ref{tab:cache-sweep} and Figure~\ref{fig:cache-sweep} report the
results across all eight thresholds.

\begin{table}[htbp]
\centering
\begin{tabular}{lccc}
\hline
\textbf{Threshold} & \textbf{Hit rate} & \textbf{Correct rate} & \textbf{False-positive rate} \\
\hline
0.80           & 100.0\% & 99.1\% & 0.9\% \\
0.85           & 100.0\% & 99.1\% & 0.9\% \\
0.88           &  99.9\% & 99.0\% & 0.9\% \\
0.90           &  99.7\% & 98.9\% & 0.9\% \\
0.92 (default) &  99.1\% & 98.3\% & 0.8\% \\
0.94           &  96.2\% & 95.4\% & 0.8\% \\
0.96           &  88.1\% & 87.5\% & 0.6\% \\
0.98           &  59.1\% & 59.1\% & 0.1\% \\
\hline
\end{tabular}
\caption{Semantic cache hit rate, correct rate, and false-positive rate over
  1974 paraphrased intents at eight similarity thresholds.}
\label{tab:cache-sweep}
\end{table}

\begin{figure}[htbp]
    \centering
    \includegraphics[width=\textwidth]{cache_sweep.png}
    \caption{Cache hit rate and correctness rate as a function of the
    similarity threshold (left axis), with the false-positive rate on the
    magnified right axis. The system default of 0.92 is marked.}
    \label{fig:cache-sweep}
\end{figure}

Between thresholds 0.80 and 0.92 the hit rate remains near 100\% and the
correct rate stays above 98\%, while the false-positive rate holds steady at
approximately 0.9\%. At the system default of 0.92, the hit rate is 99.1\%
and the correct rate is 98.3\%, with a false-positive rate of 0.8\%. Above
0.92, both the hit rate and correct rate decline sharply: at threshold 0.94
the hit rate drops to 96.2\%, at 0.96 to 88.1\%, and at 0.98 to 59.1\%.
The false-positive rate decreases in the same range but only marginally,
reaching 0.1\% at 0.98. Below 0.92, down to 0.80, the hit rate is at or
near 100\% but the false-positive rate does not decrease further, remaining at
0.9\%.

